\section{Classical Analytic Comparators}
\label{sec:analytic}

For context, we recall two standard analytic mechanisms by which
prime-supported objects fail to be represented by finite rational or
determinant data. These classical statements do not enter the proof of
the fixed-DFA prime-slice bounds.

\begin{theorem}[Prime-supported Euler products have a natural boundary]
  \label{thm:euler-product-natural-boundary}
  Let \((b_n)_{n\ge 1}\subseteq\ZZ_{\ge 0}\) and assume that
  \(b_p>0\) for infinitely many primes \(p\). Suppose that
  \[
    Z_b(z):=\prod_{n\ge 1}(1-z^n)^{-b_n}
  \]
  converges to a holomorphic function on \(\abs{z}<1\). Then the unit
  circle is a natural boundary for \(Z_b\).
\leanverified{paper_prime_languages_euler_product_natural_boundary}
\end{theorem}

\begin{proof}
  The product has integer Taylor coefficients and radius of convergence
  one. If it continued across a boundary arc, the P\'olya--Carlson
  theorem would force \(Z_b\) to be rational
  \cite{Polya1923,Carlson1921}. Its logarithmic derivative would then
  be rational. On the other hand,
  \[
    \frac{zZ_b'(z)}{Z_b(z)}
    =\sum_{m\ge1}\left(\sum_{n\mid m} n b_n\right)z^m,
  \]
  and at primes \(p\) the coefficient is \(c_1+pb_p\). Since
  \(b_p>0\) for infinitely many primes, these coefficients are
  unbounded on a prime subsequence, contradicting the eventual
  recurrence structure forced by rationality. This is the standard
  P\'olya--Carlson--Estermann argument
  \cite{Polya1923,Carlson1921,Estermann1928}.
\end{proof}

\begin{remark}
  \Cref{thm:euler-product-natural-boundary} is a
  prime-supported instance of the classical rational-versus-natural
  boundary mechanism for Euler products in the work of P\'olya, Carlson,
  and Estermann \cite{Polya1923,Carlson1921,Estermann1928}.
\end{remark}

\begin{proposition}[Periodic scalar determinant models after \(z=e^{-s}\)]
  \label{prop:finite-zeta-imaginary-periodicity}
  Let \(M\) be a finite square complex matrix. The scalar determinant
  model
  \[
    F_M(s):=\det(I-(e^{-s})M)^{-1}
  \]
  is \(2\pi i\)-periodic wherever it is holomorphic. Here \(e^{-s}\) is
  a scalar multiplying \(M\), not an exponential of the matrix. Indeed,
  \[
    F_M(s+2\pi i)
    =\det(I-(e^{-(s+2\pi i)})M)^{-1}
    =\det(I-(e^{-s})M)^{-1}
    =F_M(s).
  \]
  In particular, if \(A\) is the adjacency matrix of a shift of finite
  type \(X_A\), then its Artin--Mazur zeta function is
  \[
    \zeta_{X_A}(z)=\det(I-zA)^{-1}.
  \]
  Consequently \(\zeta_{X_A}(e^{-s})=\det(I-(e^{-s})A)^{-1}\) is
  \(2\pi i\)-periodic in \(s\).
  More generally, let
  \[
    F(s)=\sum_{n\ge 1} a_n n^{-s}
  \]
  be a Dirichlet series with abscissa of convergence \(\sigma_c\). Assume
  that there exists at least one index \(n\ge 2\) with \(a_n\neq 0\),
  and let
  \[
    n_1:=\min\{n\ge 2:a_n\neq 0\}.
  \]
  Then \(F\) is not \(2\pi i\)-periodic on any right half-plane. Moreover
  \(F\) cannot agree with \(\zeta_{X_A}(e^{-s})\) on any non-empty open
  subset of \(\Re(s)>\sigma_c\) on which the determinant model is
  holomorphic; equivalently, the two functions do not agree
  meromorphically on any open subset of that half-plane. In particular,
  the Riemann \(\zeta\)-function cannot agree with
  \(\det(I-(e^{-s})A)^{-1}\) on any non-empty open subset of \(\Re(s)>1\)
  where the determinant model is holomorphic.
\end{proposition}

\begin{proof}
  The determinant formula is classical for shifts of finite type; see
  Artin and Mazur \cite{ArtinMazur1965}, Bowen and Lanford
  \cite{BowenLanford1970}, and
  \cite[Chapter~6]{LindMarcus1995}. The scalar determinant periodicity
  has already been computed in the statement. For a Dirichlet series,
  compare \(F(\sigma+2\pi i)-F(\sigma)\) as \(\sigma\to+\infty\). The
  first non-zero term with index \(n_1\ge2\) contributes
  \[
    a_{n_1}n_1^{-\sigma}(e^{-2\pi i\log n_1}-1),
  \]
  while all later terms are \(o(n_1^{-\sigma})\) in a sufficiently far
  right half-plane. The displayed coefficient is non-zero, since
  \(\log n_1\notin\ZZ\) for an integer \(n_1\ge2\). Thus \(F\) is not
  \(2\pi i\)-periodic. Equality with a determinant model on a non-empty
  open set would propagate by the identity theorem to large real
  points in the same right half-plane component, contradicting this
  non-periodicity.
\end{proof}

\begin{remark}
  \Cref{thm:euler-product-natural-boundary,prop:finite-zeta-imaginary-periodicity}
  are complementary background tests. The first detects the boundary
  geometry of prime-supported Euler products. The second detects the
  periodicity forced on scalar determinant models by the logarithmic
  change of variables.
\end{remark}
